%==================================================================
% BAB 3 METODE PENELITIAN
% Berisi waktu dan tempat, alat dan bahan, tahapan penelitian,
% rancangan sistem, pengujian, dan jadwal pelaksanaan.
%==================================================================
\chapter{METODE PENELITIAN}
\label{chap:metode}

%------------------------------------------------------------------
\section{Waktu dan Tempat Pelaksanaan}
%------------------------------------------------------------------

Penelitian dilaksanakan pada bulan Februari sampai dengan Juni
\tahun. Pengambilan data citra dilakukan di Pasar Banyuwangi, sedangkan
proses pengolahan citra dan pelatihan jaringan saraf tiruan dilakukan
di Laboratorium Komputer Program Studi \prodi\ \institusi.

%------------------------------------------------------------------
\section{Alat dan Bahan}
%------------------------------------------------------------------

Alat dan bahan yang digunakan pada penelitian ini disajikan pada
\cref{tab:alatbahan}.

\begin{table}[H]
\caption{Alat dan Bahan Penelitian}
\label{tab:alatbahan}
\begin{spacing}{1}
\centering
\begin{tabularx}{\textwidth}{|>{\centering\arraybackslash}p{0.8cm}|X|>{\centering\arraybackslash}p{2.4cm}|}
\hline
\rowcolor{blue!12}
\textbf{No} & \textbf{Nama Alat/Bahan} & \textbf{Jumlah} \\
\hline
1 & Kamera digital 12 megapiksel & 1 unit \\
\hline
2 & Kotak akusisi citra berukuran 40 cm $\times$ 30 cm $\times$ 30 cm & 1 unit \\
\hline
3 & Lampu \textit{light emitting diode} (LED) 12 watt & 2 buah \\
\hline
4 & Ikan bandeng segar & 100 ekor \\
\hline
5 & Komputer dengan prosesor \textit{Intel Core i5} & 1 unit \\
\hline
6 & Perangkat lunak Python 3.11 dan OpenCV & 1 paket \\
\hline
\end{tabularx}
\par\vspace{2pt}
\raggedright\noindent Sumber: perencanaan penelitian (2024)
\end{spacing}
\end{table}

%------------------------------------------------------------------
\section{Tahapan Penelitian}
%------------------------------------------------------------------

Penelitian dilaksanakan melalui beberapa tahap yang saling berkaitan.
Tahapan tersebut digambarkan pada \cref{fig:tahapan}.

\begin{figure}[H]
\centering
\begin{tikzpicture}[node distance=0.5cm]
  \node (mulai)    [mulai]                        {Mulai};
  \node (studi)    [proses, below=of mulai]       {Studi literatur dan\\penyusunan instrumen};
  \node (akuisisi) [proses, below=of studi]       {Akuisisi citra mata\\dan insang ikan};
  \node (pra)      [proses, below=of akuisisi]    {Praproses: \textit{resize},\\perbaikan kontras};
  \node (ciri)     [proses, below=of pra]         {Ekstraksi ciri warna\\pada ruang HSV};
  \node (latih)    [proses, below=of ciri]        {Pelatihan jaringan\\saraf tiruan};
  \node (uji)      [proses, below=of latih]       {Pengujian model};
  \node (keputusan)[keputusan, below=of uji]      {Akurasi\\$\geq 90\%$?};
  \node (analisis) [proses, below=of keputusan]   {Analisis hasil dan\\penyusunan laporan};
  \node (selesai)  [mulai, below=of analisis]     {Selesai};

  \draw [panah] (mulai)    -- (studi);
  \draw [panah] (studi)    -- (akuisisi);
  \draw [panah] (akuisisi) -- (pra);
  \draw [panah] (pra)      -- (ciri);
  \draw [panah] (ciri)     -- (latih);
  \draw [panah] (latih)    -- (uji);
  \draw [panah] (uji)      -- (keputusan);
  \draw [panah] (keputusan) -- node[right, font=\small] {Ya} (analisis);
  \draw [panah] (analisis) -- (selesai);

  % Jalur kembali bila akurasi belum tercapai
  \draw [panah] (keputusan.east) -- ++(1.2,0)
        |- node[pos=0.25, right, font=\small] {Tidak} (latih.east);
\end{tikzpicture}
\caption{Diagram Alir Tahapan Penelitian}
\label{fig:tahapan}
\end{figure}

%------------------------------------------------------------------
\section{Perancangan Sistem}
%------------------------------------------------------------------

\subsection{Akuisisi Citra}

Pengambilan citra dilakukan di dalam kotak akusisi dengan dua lampu LED
yang dipasang pada sudut 45 derajat terhadap bidang objek. Jarak kamera
ke objek ditetapkan 25 cm dan seluruh citra diambil dengan pengaturan
ISO, kecepatan rana, serta panjang fokus yang sama agar intensitas
cahaya antar citra tetap seragam.

\subsection{Praproses Citra}

Citra hasil akuisisi disekala ulang menjadi $256 \times 256$ piksel
menggunakan interpolasi \textit{bilinear}. Perbaikan kontras dilakukan
dengan perataan histogram untuk memperjelas perbedaan warna pada mata
dan insang ikan.

\subsection{Ekstraksi Ciri Warna}

Ekstraksi ciri dilakukan pada ruang warna HSV menggunakan
\cref{eq:hsv}. Nilai rata-rata dan simpangan baku setiap kanal
dihitung dan dijadikan satu vektor ciri berukuran enam dimensi, yaitu
$\bar{H}$, $\bar{S}$, $\bar{V}$, $s_H$, $s_S$, dan $s_V$.

\subsection{Klasifikasi}

Vektor ciri dinormalisasi ke rentang 0 sampai 1, kemudian digunakan
sebagai masukan jaringan saraf tiruan dengan arsitektur 6 neuron
masukan, dua lapisan tersembunyi berisi 16 dan 8 neuron, serta 3 neuron
keluaran yang mewakili kelas segar, cukup segar, dan tidak segar.

%------------------------------------------------------------------
\section{Pengujian dan Evaluasi}
%------------------------------------------------------------------

Pengujian dilakukan menggunakan \textit{k-fold cross validation} dengan
$k = 5$. Kinerja sistem diukur menggunakan akurasi seperti pada
\cref{eq:akurasi}.

\begin{equation}
\mathrm{Akurasi} = \frac{\text{jumlah data uji yang benar}}{\text{jumlah seluruh data uji}} \times 100\%
\label{eq:akurasi}
\end{equation}

\noindent Selain akurasi, dihitung pula presisi dan perolehan
(\textit{recall}) untuk setiap kelas menggunakan
\cref{eq:presisi,eq:recall}.

\begin{equation}
\mathrm{Presisi} = \frac{TP}{TP+FP} \times 100\%
\label{eq:presisi}
\end{equation}

\begin{equation}
\mathrm{Perolehan} = \frac{TP}{TP+FN} \times 100\%
\label{eq:recall}
\end{equation}

\noindent dengan $TP$ adalah \textit{true positive}, $FP$ adalah
\textit{false positive}, dan $FN$ adalah \textit{false negative}.

%------------------------------------------------------------------
\section{Jadwal Pelaksanaan}
%------------------------------------------------------------------

Jadwal pelaksanaan penelitian disusun dalam bentuk \textit{bar-chart}
seperti pada \cref{tab:jadwal}.

\begin{table}[H]
\caption{Model \textit{Bar-Chart} Penyusunan Jadwal Pelaksanaan Tugas Akhir}
\label{tab:jadwal}
\begin{spacing}{1}
\centering
\begin{tabular}{|>{\raggedright\arraybackslash}p{5.2cm}|*{6}{>{\centering\arraybackslash}p{0.9cm}|}}
\hline
\multirow{2}{*}{\textbf{Jenis Kegiatan}} & \multicolumn{6}{c|}{\textbf{Bulan Ke-}} \\
\cline{2-7}
 & \textbf{1} & \textbf{2} & \textbf{3} & \textbf{4} & \textbf{5} & \textbf{6} \\
\hline
Studi literatur            & \cellcolor{blue!25} & & & & & \\
\hline
Penyusunan instrumen       & \cellcolor{blue!25} & \cellcolor{blue!25} & & & & \\
\hline
Akuisisi citra             & & \cellcolor{blue!25} & & & & \\
\hline
Praproses dan ekstraksi    & & & \cellcolor{blue!25} & & & \\
\hline
Pelatihan dan pengujian    & & & & \cellcolor{blue!25} & & \\
\hline
Analisis dan penulisan     & & & & & \cellcolor{blue!25} & \cellcolor{blue!25} \\
\hline
\end{tabular}
\par\vspace{2pt}
\raggedright\noindent Sumber: perencanaan penelitian (2024)
\end{spacing}
\end{table}
